\documentclass{article}
\usepackage[a4paper]{geometry}
\usepackage{fancyhdr}
\pagestyle{fancy}
\lhead{Ethisches Bewerten}
\rhead{März 2026}
\begin{document}
\section{Ethisches Bewerten}
Soll ein Sachverhalt oder Ähnliches ethisch bewertet werden, so ist sechs Schritten zu folgen
\begin{enumerate}
 \item Das vorliegende ethische Problem und dessen sachlichen Grundlagen in eigenen Worten erklären.
 \item Die Lösungen, in der Regel Handlungsoptionen einer Person, aufzählen.
 \item Die Pro- und Kontraargumente aller Möglichkeiten aufzählen.
 \item Die ethische Werte, die hinter den Argumenten stecken, aufzählen. Hier muss zwischen \emph{deontologischen} und \emph{konsequenzialistischen} Argumentationsweisen unterschieden werden.
 \item Das eigene, persönliche, Urteil begründet fällen und andere Urteile diskutieren.
 \item Konsequenzen des eigenen und der anderen Urteile aufzählen.
\end{enumerate} 
 
\subsection{Argumentationsweisen}
\begin{description}
 \item[konsequenzialistisch] wird das die Konsequenz mitbetrachtet. Etwas woraus gutes folgt ist gut.
 \item[deontologisch] Betrachtet die Tat unabhängig von ihrer Konsequenz. Etwas schlechte zu tun, auch wenn gutes folgt, ist schlecht.
\end{description} 
\end{document}